% ============================================================
%  Section 3 — Background  (~0.75 pages)
% ============================================================
\section{Background}

\subsection{Frameworks}

\textbf{Qiskit}~\cite{qiskit} is IBM's open-source SDK built around the
\texttt{QuantumCircuit} abstraction.
Its multi-level transpiler pipeline supports four optimisation passes
(levels 0--3) that progressively reduce gate count and circuit depth
through algebraic simplification and routing heuristics.

\textbf{Cirq}~\cite{cirq} is Google's hardware-oriented SDK designed
around a grid-qubit model that reflects physical device topology.
Circuits are organised into \emph{moments}---layers of simultaneously
executable gates---giving the programmer explicit control over
parallelism and depth.

\textbf{PennyLane}~\cite{bergholm2018pennylane} is Xanadu's
differentiable programming SDK whose central abstraction is the
device-agnostic \texttt{QNode}.
Its primary design goal is variational and machine-learning workflows,
where automatic differentiation of quantum circuits is essential.
Gate decompositions are chosen for portability across heterogeneous
backends rather than for minimising gate count on a specific target.

\subsection{OpenQASM 3.0}
OpenQASM 3.0~\cite{cross2022openqasm3} is the lingua franca that enables
the unified pipeline in this work.
Over its predecessor (QASM 2.0), version 3.0 adds: reusable subroutines
and gate definitions with classical parameters; complex numeric types;
gate modifiers (\texttt{ctrl@}, \texttt{pow@}, \texttt{inv@}); and
structured classical control flow (\texttt{if}, \texttt{for},
\texttt{while}).
These extensions allow QCanvas to faithfully represent circuits from all
three frameworks in a single, backend-agnostic intermediate
representation without loss of gate semantics.

\subsection{Metrics}

\textbf{Jensen--Shannon Divergence (JSD)}~\cite{lin1991divergence} is a
symmetric, bounded measure of distributional similarity on $[0, 1]$.
Unlike KL divergence, JSD is defined even when the two distributions
have non-overlapping support, making it suitable for comparing
measurement outcome distributions that may differ in their set of
reachable basis states.
We declare two distributions statistically equivalent when
$\text{JSD} < 0.02$.

\textbf{Quantum Volume}~\cite{cross2019validating} is a single-number
circuit quality estimator equal to $2^n$ where $n$ is the largest number
of qubits for which a model random circuit achieves heavy-output
generation probability above $2/3$.
We compute QV in a structure-implied mode from circuit depth and width
under a depolarizing noise model~\cite{nielsen2010quantum}.

\textbf{Power-law scaling} fits the relation $y = c \cdot n^a$ to gate
count or depth as qubit count $n$ varies.
The exponent $a$ classifies the scaling regime: $a \approx 0$ is
constant (optimal for oracle algorithms), $a = 1$ is linear, and
$a > 1$ is super-linear.
